%% One-page ATS-friendly resume template
%% ----------------------------------------------------------------------------
%% Single-column, icon section headers, dot-rated languages, pipe-separated
%% skills. Modeled on a clean modern resume; profile-agnostic ([PLACEHOLDER]
%% tokens) so it is safe to commit and share.
%%
%% Compile with: lualatex -interaction=nonstopmode <file>.tex
%%   - Uses the `lato` package (no font-path juggling; compiles from any dir).
%%   - Uses fontawesome5 for the section/contact icons.
%%   - pdflatex can fail on modern MiKTeX with fontawesome5 font-expansion
%%     errors; lualatex handles it cleanly.
%%
%% ATS notes (why this layout is parser-safe):
%%   - Single column => extracted reading order matches the visual order.
%%   - Every contact detail (email, phone) is printed as LITERAL TEXT next to
%%     its icon, never carried by the icon alone.
%%   - The contact block is ONE line (email, phone, location, LinkedIn,
%%     Portfolio) so the page keeps that second line for real content. The
%%     Portfolio item is the one deliberate exception to the literal-text rule:
%%     its label is the word "Portfolio" and the URL lives only in the link
%%     target, which is what Salman asked for.
%%   - Lato embeds a proper Unicode map => no (cid:*) garbage on extraction.
%%   - Section headings are real text; the icons are decoration only.
%% ----------------------------------------------------------------------------

\documentclass[10pt,a4paper]{article}

% --- Font: Lato as the default family (sans, clean, ATS-safe) ---
\usepackage[T1]{fontenc}
\usepackage[default]{lato}
\renewcommand{\familydefault}{\sfdefault}

% --- Layout packages ---
\usepackage[left=1.25cm,right=1.25cm,top=0.8cm,bottom=0.8cm]{geometry}
\usepackage{fontawesome5}
\usepackage{enumitem}
\usepackage{xcolor}
\usepackage{tabularx}
\usepackage{array}
\usepackage{needspace}
\usepackage{ifthen}
\usepackage[hidelinks]{hyperref}

% --- Accent colour (affects rendering only; irrelevant to text extraction) ---
\definecolor{accent}{HTML}{1F3B57}   % dark slate blue
\definecolor{dim}{HTML}{555555}      % muted grey for secondary text
\definecolor{dotoff}{HTML}{C8CDD3}   % empty language-rating dot

% --- Global spacing / no page numbers ---
\pagestyle{empty}
\setlength{\parindent}{0pt}
\setlength{\tabcolsep}{0pt}
\linespread{0.92}   % tight but readable; adjust 0.90–1.00 to land on exactly one page

% --- Hyperlink metadata ---
\hypersetup{
  pdftitle={[YOUR_NAME] - Resume},
  pdfauthor={[YOUR_NAME]},
}

% ============================================================================
%  MACROS
% ============================================================================

% Section header: icon + bold coloured title + full-width rule.
%   \cvsection{\faIcon}{TITLE}
\newcommand{\cvsection}[2]{%
  \needspace{2\baselineskip}%
  \vspace{2pt}%
  {\color{accent}\large\bfseries #1\hspace{6pt}#2}\par
  \vspace{-5pt}{\color{accent}\rule{\linewidth}{0.9pt}}\par
  \vspace{2pt}%
}

% Experience / education entry header.
%   \cventry{Role}{Organisation}{Date range}{Location}
% "Role, Organisation" (bold role + italic org) at left, muted "date | location"
% flush right on the SAME line. One visual line per entry instead of a
% three-line stack, which is where a one-pager wins back vertical space for
% actual content (~2 lines per entry).
% \parfillskip=0pt makes this final line fully justified so \hfill absorbs all
% the slack; without it \parfillskip steals half and the date drifts inward.
\newcommand{\cventry}[4]{%
  \needspace{3\baselineskip}%
  {\parfillskip=0pt\relax
   \textbf{#1,}\ {\itshape\color{accent}#2}\hfill
   {\color{dim}#3 \textbar{} #4}\par}%
  \vspace{1pt}%
}

% Tight bullet list for achievements. Bold the metrics/keywords inside.
% Target ~1 \textbf{} call per bullet (experience + project bullets) — more
% than that crowds the page and weakens emphasis. Skill-group labels are
% exempt; they are always bold by design.
\newenvironment{cvbullets}{%
  \begin{itemize}[leftmargin=1.35em,itemsep=0.5pt,parsep=0pt,topsep=0.5pt,%
                  label={\color{accent}\small\textbullet}]
}{%
  \end{itemize}%
}

% A labelled skills line: \skillgroup{Category}{item | item | item}
\newcommand{\skillgroup}[2]{%
  \textbf{#1:}~ #2\par\vspace{1.5pt}%
}

% Language 5-dot rating: \rating{n} draws n filled + (5-n) empty dots.
\newcounter{ratingdot}
\newcommand{\rating}[1]{%
  \setcounter{ratingdot}{0}%
  \whiledo{\value{ratingdot} < #1}{\textcolor{accent}{\scriptsize\faCircle}\,\stepcounter{ratingdot}}%
  \whiledo{\value{ratingdot} < 5}{\textcolor{dotoff}{\scriptsize\faCircle}\,\stepcounter{ratingdot}}%
}

% One language row: \lang{Name}{filled-dots}{Text level}
\newcommand{\lang}[3]{%
  \makebox[3.6cm][l]{#1}\makebox[3.2cm][l]{\rating{#2}}{\color{dim}#3}\par\vspace{1pt}%
}

% Separator between contact items. Deliberately tighter than \quad (1em = 9pt at
% \small) so all five items fit on ONE line inside \textwidth. The contact block
% is a single line by contract — the second line it used to occupy is reclaimed
% for experience/projects content. If an item is added and the line overflows,
% shorten an item's text rather than widening the margins.
\newcommand{\contactsep}{\hspace{7pt}}

% ============================================================================
%  DOCUMENT
% ============================================================================
\begin{document}

% ---------- HEADER ----------
{\centering
  {\color{accent}\fontsize{24}{28}\selectfont\bfseries [YOUR_NAME]}\par
  \vspace{2pt}
  {\color{dim}\normalsize [YOUR_TAGLINE]}\par   % e.g. "AI/ML Engineer | Data Science"
  \vspace{3pt}
  {\small
    \faEnvelope\ \href{mailto:[YOUR_EMAIL]}{[YOUR_EMAIL]}\contactsep
    \faPhone*\ [YOUR_PHONE]\contactsep
    \faMapMarker*\ [YOUR_LOCATION]\contactsep
    \faLinkedin\ \href{[YOUR_LINKEDIN_URL]}{[YOUR_LINKEDIN_HANDLE]}\contactsep
    \faGlobe\ \href{[YOUR_PORTFOLIO_URL]}{Portfolio}%
  }\par
}
\vspace{2pt}

% ---------- SUMMARY ----------
\cvsection{\faUserTie}{SUMMARY}
[YOUR_PROFILE_STATEMENT — 2-3 lines. Tailor per role: lead with the single
strongest match to the posting, then breadth. No first-person pronoun.]

% ---------- EXPERIENCE ----------
\cvsection{\faBriefcase}{PROFESSIONAL EXPERIENCE}

\cventry{[ROLE_TITLE]}{[COMPANY]}{[START] -- [END]}{[CITY, COUNTRY]}
\begin{cvbullets}
  \item [Achievement with a single \textbf{bolded metric/keyword} and an action verb.]
  \item [Achievement mapped to a posting requirement.]
  \item [Achievement showing scope / tools.]
\end{cvbullets}
\vspace{3pt}

\cventry{[ROLE_TITLE]}{[COMPANY]}{[START] -- [END]}{[CITY, COUNTRY]}
\begin{cvbullets}
  \item [Achievement.]
  \item [Achievement.]
\end{cvbullets}
% ... repeat \cventry + cvbullets per role. Relevance-weighted cutting keeps
% the whole document to one page (see 05-cv-templates.md).
% BOLD RULE: aim for ~1 \textbf{} call per experience/project bullet — more
% crowds the page and weakens emphasis. Skill-group labels are always bold.

% ---------- SKILLS ----------
\cvsection{\faTools}{SKILLS}
\skillgroup{[CATEGORY_1]}{[item | item | item | item]}
\skillgroup{[CATEGORY_2]}{[item | item | item | item]}
\skillgroup{[CATEGORY_3]}{[item | item | item | item]}

% ---------- PROJECTS ----------
\cvsection{\faLightbulb}{PROJECTS}
\begin{cvbullets}
  \item \textbf{[PROJECT_NAME]} ([tech stack]) -- [one-line outcome].
  \item \textbf{[PROJECT_NAME]} ([tech stack]) -- [one-line outcome].
\end{cvbullets}

% ---------- EDUCATION ----------
\cvsection{\faGraduationCap}{EDUCATION}
\cventry{[DEGREE]}{[INSTITUTION]}{[START] -- [END]}{[CITY, COUNTRY]}
{\color{dim}[One line: thesis / focus / notable coursework.]}\par
\vspace{3pt}

% ---------- LANGUAGES ----------
\cvsection{\faLanguage}{LANGUAGES}
\lang{[LANGUAGE]}{5}{[Native / bilingual]}
\lang{[LANGUAGE]}{4}{[Professional working proficiency]}
\lang{[LANGUAGE]}{2}{[Elementary (A2)]}

% Certifications can be folded into EDUCATION via \skillgroup{Certifications}{...}
% (see Sample CV/main.tex for the pattern) rather than using a separate AWARDS section.

\end{document}
